% No preamble. Prefix mcr-. This is a continuum reference, not a lattice comparison.
\section{A completely specified massive continuum reference}
\label{mcr-section}

We compute a reference with its regulator and normalization fixed.
This supplies a target for comparison, not an identification of
the actual constrained lattice determinant with that target.
The background-field operators and heat coefficient agree with
the conventions described in \cite{vassilevichctn}, Sections 3.4,
4.2.1 and 8.3. The calculation below keeps the regulator and
the one-color Maxwell extraction explicit.

Use final physical coordinates on \(\mathbb R^4\), scalar mass
\(m>0\), and a lower proper-time cutoff \(\epsilon^2>0\).
Choose a Hermitian Cartan element \(H\) with
\(\tau(H^2)=1\). Let \(q\) run over the eigenvalues of
\(\operatorname{ad}_H\) on the complexified adjoint space,
with multiplicities. Both signs of a nonzero charge are
included. Complexification has the same trace as the
original real color space; it does not introduce another
factor of two. Put
\[
 C_H=\sum_q q^2.
\]
For one constant curvature component \(F_{12}=fH\), a charge
sector has scalar operator \(L_q=-D_\mu D_\mu\), where one
may choose \(D_2=\partial_2+iqf x_1\) and the other derivatives
ordinary. The vector reference is
\[
 (V_q)_{\mu\nu}=L_q\delta_{\mu\nu}-2iqF^{(1)}_{\mu\nu},
 \qquad F^{(1)}_{12}=f,\quad F^{(1)}_{21}=-f.
\]
The spin matrix has eigenvalues \(0,0,2qf,-2qf\).
Its sign convention has no effect on its trace.
We add \(m^2\) to both operators.

\subsection{Exact diagonal heat traces}
\label{mcr-heat-section}

Fourier transformation in \(x_2\) reduces the magnetic
plane to oscillators with eigenvalues \(|qf|(2n+1)\).
For a normalized oscillator eigenfunction, integration
of its shifted squared value gives
\[
 \int_{\mathbb R}\frac{dp_2}{2\pi}
  \left|\psi_n\!\left(x_1+\frac{p_2}{qf}\right)\right|^2
       =\frac{|qf|}{2\pi}.
\]
Thus this is a pointwise diagonal density calculation,
without a finite-box degeneracy assumption. Summing the
oscillator levels and multiplying by the heat density
of the two free coordinates gives
\begin{equation}
 k_q(s;f)=(4\pi s)^{-2}\frac{qfs}{\sinh(qfs)}.
 \label{mcr-scalar-heat}
\end{equation}
At \(qf=0\) the fraction is defined by its limit.
The vector density is
\[
 k_q(s;f)\{2+2\cosh(2qfs)\}.
\]

Define the normalized proper-time determinant correction
with weights \(-\tfrac12\log\det(V+m^2)+\log\det(L+m^2)\).
Here this phrase means the following density integral;
there is no infinite-volume determinant:
\begin{equation}
 {\cal R}_{\epsilon,m}(f)
 =\frac1{16\pi^2}\int_{\epsilon^2}^{\infty}
       \frac{ds}{s^3}e^{-m^2s}
       \sum_q\left\{\frac{qfs}{\sinh(qfs)}
                          \cosh(2qfs)-1\right\}.
 \label{mcr-reference}
\end{equation}
The sign follows from
\(\log P=-\int ds\,e^{-sP}/s\), after subtraction and
the specified proper-time cutoff. The two vector
polarizations with zero spin cancel the scalar ghost
in this combined expression.

For real \(|f|\max_q|q|<m^2\), the massive vector has
positive bottom and the integral converges at infinity.
On a smaller real interval, differentiation twice at
zero is justified by a common exponentially decaying
majorant. In particular, the notation \([f^2]\) below
denotes half the second real derivative at zero.
We do not assert a convergent complex Taylor series
in a nonzero constant magnetic field.

\begin{proposition}[Normalized reference coefficient]
\label{mcr-coefficient}
Set \(E_1(x)=\int_x^\infty e^{-u}\,du/u\). In the
normalization \(\beta=2[f^2]\) used in this notebook,
\begin{equation}
 \beta_{\rm ref}(\epsilon,m)
       =\frac{11C_H}{48\pi^2}E_1(m^2\epsilon^2).
 \label{mcr-beta}
\end{equation}
In particular, with \(m=1\), \(\epsilon=K^{-1}\),
\begin{equation}
 \left|\beta_{\rm ref}(K^{-1},1)
           -\frac{11C_H}{24\pi^2}\log K\right|
 \le \frac{11C_H}{48\pi^2}(1+e^{-1}),\qquad K\ge1.
 \label{mcr-finite-part}
\end{equation}
\end{proposition}

\begin{proof}
The scalar magnetic factor and the spin factor give
\[
 \frac{x}{\sinh x}\cosh(2x)
       =1+\frac{11}{6}x^2+O(x^4).
\]
Differentiating \eqref{mcr-reference} twice at zero
therefore gives
\[
 [f^2]{\cal R}_{\epsilon,m}
    =\frac{11C_H}{96\pi^2}
           \int_{\epsilon^2}^{\infty}\frac{e^{-m^2s}}s\,ds.
\]
The substitution \(u=m^2s\) proves \eqref{mcr-beta}.
For \(0<x\le1\),
\[
 E_1(x)-\log(1/x)
   =\int_x^1\frac{e^{-u}-1}{u}\,du
                  +\int_1^\infty\frac{e^{-u}}u\,du.
\]
The absolute value is at most \(1+e^{-1}\). This proves
the asserted uniform finite part.
\end{proof}

For the repository's \(SU(N)\) normalization
\(\tau=\operatorname{Tr}_{\rm fund}/N\), write
\(H=\operatorname{diag}(h_1,\ldots,h_N)\). Then
\(\sum_i h_i=0\), \(\sum_i h_i^2=N\), and the
nonzero-color charges are \(h_i-h_j\), \(i\ne j\);
the \(N-1\) Cartan charges vanish. Consequently
\[
 C_H=\sum_{i,j}(h_i-h_j)^2
              =2N\sum_i h_i^2=2N^2.
\]
The coefficient of \(\log K\) is thus
\(11N^2/(12\pi^2)\), or \(11/(3\pi^2)\) for \(SU(2)\),
in this one-normalized-color convention.
This is not a sum of \(N^2-1\) different Maxwell extractions.

\subsection{Scale and regulator bookkeeping}
\label{mcr-scaling-section}

The same continuum reference can be written in current
fine units, with cutoff \(s\ge1\), mass \(K^{-1}\), and
curvature \(f/K^2\). Multiplication of that density by
the block volume \(K^4\), followed by \(s=K^2u\), gives
exactly \eqref{mcr-reference} with \(\epsilon=K^{-1}\),
\(m=1\). Thus there is no additional factor of \(K^4\)
in its extracted coefficient. This is a change of units
within the continuum reference, not a replacement of a
lattice heat kernel by a continuum kernel.

Fixed positive changes of these continuum regulators
have a bounded effect. For \(c,m>0\),
\begin{equation}
 |\beta_{\rm ref}(c/K,m)-\beta_{\rm ref}(1/K,1)|
 \le\frac{11C_H}{24\pi^2}|\log(cm)|.
 \label{mcr-regulator-change}
\end{equation}
Indeed, the difference of the two \(E_1\) integrals is
bounded by the integral of \(du/u\) between their lower
limits. Uniformity follows when \(cm\) remains in a
fixed compact subinterval of \((0,\infty)\).
This estimate does not compare a based ghost to an
unbased massive ghost, change a lattice constraint, or
control a massless zero-mode limit.

